\begin{proposition}[Proved: the complete star cardinality reduction]
\label{CP2:A-star-seventeen-scope}
To exclude all three-object categories with star-shaped acyclic
support and at most seventeen arrows, it suffices to exclude
the $38$ exact seventeen-arrow profiles in
\path{converse/scope17/star17_profiles.tsv}.
Each has center order between five and seven, and both leaf
orders at most seven. No profile outside this explicit list
requires a composition-table search.
\end{proposition}
\begin{proof}
Identity padding, Proposition~\ref{CP2:A-identity-padding},
reduces to exactly seventeen arrows without changing the
support or classifying space. Taking opposition makes the
center a source. A profile is a positive tuple
$(a,b,c,s,t)$ with sum seventeen, in which $(b,s)$ and
$(c,t)$ may be interchanged. There are
$\binom{16}{4}=1820$ labelled tuples. The leaf-swap fixed
tuples satisfy $a+2b+2s=17$; their number is
$\sum_{d=2}^{8}(d-1)=28$. Thus there are exactly
$(1820+28)/2=924$ orbits. The fresh program
\path{converse/scope17/star17_scope.py} retains the representative
with $(b,s)\le(c,t)$ in lexicographic order and tests only
the following proved sufficient conditions.

The large-leaf, large-center and small-center results above
remove every corner configuration outside center orders
five through seven and leaf orders at most seven.
For completeness the program also records the overlapping
elementary cuts: all three corner second cohomology groups
zero; a center-fixed singleton stalk when both leaves have
zero second cohomology; a full edge with zero second
cohomology and a remaining leaf with the same property;
the coned-binary lemma; and the small-corner binary-stalk
lemma. The sufficient edge tests for zero second cohomology
are a singleton stalk, total edge size at most five, or
both corner orders at most two. They follow respectively
from the join model, the second-homology cost lemma, and
semisimplicity of the two corner algebras.

There are two further transfer tests. If the center order
is at most five, its rational homology above degree two
vanishes. The star is the union of its two full edges
along that corner. Mayer--Vietoris transfers boundedness
to both edges and makes restriction in fourth cohomology
injective. It is therefore enough that both edges have
zero second products. The program uses only: stalk size
at most two; both corner orders at most four; total edge
size at most eleven; a corner of order at most two opposite
a corner of order at most nine; or a corner of order at
most three opposite a corner of order at most six.
The first three are the binary, two-block and small-edge
lemmas. The fourth is exact orbit comparison followed by
the bounded-nine act theorem. For the fifth, the only
nonprojective three-element endpoint is the relevant
orientation of the identity-adjoined right-zero band;
its natural graph model and
Corollary~\ref{CP2:A-baseline-right-zero-three-six} apply.
All uses of boundedness occur after this Mayer--Vietoris
transfer.

Finally, if both leaves have orders at most four, each
nonprojective leaf is allowed only when its stalk has at
most two states. Lemma~\ref{CP2:A-star-leaf-reductions}
then either gives a center-fixed state or reduces both
leaves to trivial corners. Merging the leaves gives a
bounded ordinary act cone, excluded by the bounded-nine
theorem when the center order is at most nine.

The exact enumeration leaves $38$ profiles and removes
$886$. The full reason ledger, including every overlapping
predicate, is archived as
\path{outputs/star17/cardinality_scope.json}; its source
and ledger digests are in
\path{converse/scope17/star17_scope_summary.json}.
The maximum number of normalized two-simplices among the
residual profiles is $88$. The cardinality count and all
reduction predicates received a separate hostile proof
check. This proposition establishes the complete scope
reduction; it does not treat an unfinished composition
stream as a completed exclusion.
\end{proof}

\paragraph{Composition enumeration for the residual stars.}
The executable \path{converse/scope17/star17_canonical.cpp}
uses the action-orbit procedure proved in
Proposition~\ref{CP2:A-linear-biset-enumeration}. For a source
star there are exactly two independent connecting bisets,
and no composition of two cross arrows. Thus a pair of
canonical bisets supplies the complete category table.
If their state stabilizer orders are $A_S,A_T$, the
weight $s!t!/(A_SA_T)$ restores every labelled table.
All ordinary corner tables through order seven are used.

Only three types of corner-table cut precede this
enumeration. A projective trivial left module at the
center permits an actual orbit quotient to a trivial
center; the result is a wedge of bounded ordinary leaf
act cones of order at most seven. Projective trivial
right modules at both leaves permit their orbit quotients
to trivial leaves; merging them gives a bounded ordinary
center-act cone of order at most seven. The bounded-nine
theorem excludes both cases. The remaining cuts require
zero second cohomology on both blocks of a two-block
decomposition. Corner second-homology vanishing is checked
by a literal binary bar calculation, a sufficient condition
for rational vanishing by the universal coefficient theorem.
An edge is declared second-homology zero only for a singleton
stalk, two semisimple corner algebras, or total size at most
five. These tests received a separate targeted proof check.

The unchanged literal-certificate procedure checks the
necessary binary first and second homology pattern, then
produces either a square primitive or a positive third
homology calculation. The independent checker
\path{converse/scope17/check_star17.cpp} rechecks every
retained typed associative table and its bar certificate.
Only a closed, hashed stream with matching producer and
checker completion counts can certify its profile.
